% Generated from 2/2-6.md; technical content preserved.
\exercise{۲.۶}
صورت \textbf{\lr{Exercise 2.6}} در کتاب درسی به شرح زیر است:

\begin{quote}
\textbf{صورت مسئله:}

رابطه (2.27) کتاب را اثبات کنید:

\[V_{dsat} = \frac{V_{GT} V_c}{V_{GT} + V_c} = V_{GT} \parallel V_c\]

\end{quote}
این رابطه، ولتاژ اشباع درین (\(V_{dsat}\)) را با در نظر گرفتن هم‌زمان پدیده \textbf{\lr{Velocity Saturation}} (اشباع سرعت حامل‌ها) و اثر مرتبه اول ترانزیستور (\lr{Shockley Pinch-off}) توصیف می‌کند.

\subsection*{۱. تعاریف و معادلات پایه (\lr{Givens \& Governing Equations})}

\begin{enumerate}
\item \textbf{مدل سرعت حامل‌ها (\lr{Carrier Drift Velocity Model}):}
با افزایش میدان الکتریکی جانبی (\(E = V_{ds}/L\))، سرعت حرکت الکترون‌ها به دلیل برخورد با فونون‌های شبکه دچار اشباع می‌شود. مدل تقریب سرعت طبق رابطه (2.24) کتاب عبارت است از:
\[v = \frac{\mu_{eff} E}{1 + \frac{E}{E_c}} = \frac{\mu_{eff} \frac{V_{ds}}{L}}{1 + \frac{V_{ds}}{E_c L}} = \frac{\mu_{eff} \frac{V_{ds}}{L}}{1 + \frac{V_{ds}}{V_c}}\]
که در آن \(E_c\) میدان الکتریکی بحرانی (\lr{Critical Electric Field}) و \(V_c = E_c L\) ولتاژ بحرانی (\lr{Critical Voltage}) است. با توجه به شرط پیوستگی سرعت در \(E = E_c\) که سرعت به \(v_{sat}\) می‌رسد، داریم:
\[E_c = \frac{2 v_{sat}}{\mu_{eff}} \implies V_c = E_c L = \frac{2 v_{sat} L}{\mu_{eff}} \implies v_{sat} = \frac{\mu_{eff} V_c}{2L}\]

\item \textbf{جریان در ناحیه خطی (\lr{Linear Region Current}):}
با جایگذاری رابطه سرعت فوق در انتگرال بار کانال، رابطه جریان در ناحیه خطی به‌صورت زیر اصلاح می‌شود:
\[I_{ds} = \frac{\mu_{eff}}{1 + \frac{V_{ds}}{V_c}} C_{ox} \frac{W}{L} \left( V_{GT} - \frac{V_{ds}}{2} \right) V_{ds} \quad (V_{ds} < V_{dsat})\]
که در آن \(V_{GT} = V_{gs} - V_t\) معرف \lr{Overdrive Voltage} گیت است.

\item \textbf{جریان در نقطه اشباع (\lr{Saturation Region Current}):}
در نقطه شروع اشباع (\(V_{ds} = V_{dsat}\))، سرعت حامل‌ها در لبه درین به \(v_{sat}\) می‌رسد و بار وارونگی درین برابر \(Q_{channel} = C_{ox} W (V_{GT} - V_{dsat})\) است. در نتیجه جریان اشباع برابر است با:
\[I_{dsat} = C_{ox} W (V_{GT} - V_{dsat}) v_{sat} \quad (V_{ds} \ge V_{dsat})\]

\end{enumerate}

\subsection*{۲. اثبات گام‌به‌گام (\lr{Step-by-Step Derivation})}

طبق تعریف فیزیکی و پیوستگی مشخصه جریان-ولتاژ (\lr{I-V})، در مرز ناحیه خطی و اشباع یعنی در \(V_{ds} = V_{dsat}\)، جریان رابطه ناحیه خطی باید با جریان رابطه ناحیه اشباع برابر باشد:

\[\left. I_{ds,linear} \right|_{V_{ds} = V_{dsat}} = I_{dsat}\]

با قرار دادن فرمول‌های طرفین:
\[\frac{\mu_{eff}}{1 + \frac{V_{dsat}}{V_c}} C_{ox} \frac{W}{L} \left( V_{GT} - \frac{V_{dsat}}{2} \right) V_{dsat} = C_{ox} W (V_{GT} - V_{dsat}) v_{sat}\]

\subsubsection*{گام اول: ساده‌سازی جملات مشترک}
عبارت ثابت \(C_{ox} W\) را از طرفین معادله حذف می‌کنیم:
\[\frac{\mu_{eff}}{L} \frac{1}{1 + \frac{V_{dsat}}{V_c}} \left( V_{GT} - \frac{V_{dsat}}{2} \right) V_{dsat} = (V_{GT} - V_{dsat}) v_{sat}\]

\subsubsection*{گام دوم: جایگذاری سرعت اشباع (\(v_{sat}\))}
با استفاده از رابطه \(v_{sat} = \frac{\mu_{eff} V_c}{2L}\) در سمت راست:
\[\frac{\mu_{eff}}{L} \frac{1}{1 + \frac{V_{dsat}}{V_c}} \left( V_{GT} - \frac{V_{dsat}}{2} \right) V_{dsat} = (V_{GT} - V_{dsat}) \left( \frac{\mu_{eff} V_c}{2L} \right)\]

ضریب مشترک \(\frac{\mu_{eff}}{L}\) از دو طرف معادله ساده می‌شود:
\[\frac{1}{1 + \frac{V_{dsat}}{V_c}} \left( V_{GT} - \frac{V_{dsat}}{2} \right) V_{dsat} = (V_{GT} - V_{dsat}) \frac{V_c}{2}\]

\subsubsection*{گام سوم: حذف مخرج و ساده‌سازی جبری}
مخرج سمت چپ را به‌صورت یک کسر مشترک می‌نویسیم:
\[1 + \frac{V_{dsat}}{V_c} = \frac{V_c + V_{dsat}}{V_c}\]

با جایگذاری در معادله:
\[\frac{V_c}{V_c + V_{dsat}} \left( V_{GT} - \frac{V_{dsat}}{2} \right) V_{dsat} = (V_{GT} - V_{dsat}) \frac{V_c}{2}\]

با حذف \(V_c\) از دو طرف و ضرب طرفین در \(2 (V_c + V_{dsat})\):
\[2 \left( V_{GT} - \frac{V_{dsat}}{2} \right) V_{dsat} = (V_{GT} - V_{dsat})(V_c + V_{dsat})\]

\subsubsection*{گام چهارم: بسط دو طرف معادله}
سمت چپ معادله:
\[2 V_{GT} V_{dsat} - V_{dsat}^2\]

سمت راست معادله:
\[(V_{GT} - V_{dsat})(V_c + V_{dsat}) = V_{GT} V_c + V_{GT} V_{dsat} - V_c V_{dsat} - V_{dsat}^2\]

معادله حاصل به‌صورت زیر درمی‌آید:
\[2 V_{GT} V_{dsat} - V_{dsat}^2 = V_{GT} V_c + V_{GT} V_{dsat} - V_c V_{dsat} - V_{dsat}^2\]

جمله درجه دوم \(- V_{dsat}^2\) از دو طرف حذف می‌شود:
\[2 V_{GT} V_{dsat} = V_{GT} V_c + V_{GT} V_{dsat} - V_c V_{dsat}\]

\subsubsection*{گام پنجم: جداسازی \(V_{dsat}\)}
با انتقال \(V_{GT} V_{dsat}\) و \(-V_c V_{dsat}\) به سمت چپ:
\[2 V_{GT} V_{dsat} - V_{GT} V_{dsat} + V_c V_{dsat} = V_{GT} V_c\]
\[V_{GT} V_{dsat} + V_c V_{dsat} = V_{GT} V_c\]

فاکتورگیری از \(V_{dsat}\):
\[V_{dsat} (V_{GT} + V_c) = V_{GT} V_c\]

در نهایت:
\[V_{dsat} = \frac{V_{GT} V_c}{V_{GT} + V_c}\]

\subsection*{۳. تحلیل فیزیکی و رفتارهای حدی (\lr{Limiting Cases})}

این فرمول ساختاری کاملاً مشابه با فرمول مقاومت‌های موازی دارد (\(V_{dsat} = V_{GT} \parallel V_c\)):

\begin{enumerate}
\item \textbf{حالت کانال بلند (\lr{Long-Channel Regime}):}
اگر طول کانال بزرگ باشد یا میدان ضعیف باشد (\(V_c \gg V_{GT}\))، اشباع سرعت رخ نمی‌دهد. در این حالت:
\[V_{dsat} \approx V_{GT} = V_{gs} - V_t\]
که همان ولتاژ اشباع مرسوم \lr{Shockley} در اثر \lr{Pinch-off} است.

\item \textbf{حالت اشباع سرعت شدید (\lr{Velocity-Saturated Regime}):}
اگر طول کانال بسیار کوتاه باشد یا \(V_c \ll V_{GT}\)، حامل‌ها خیلی سریع پیش از پینچ‌آف شدن به حداکثر سرعت ممکن می‌رسند:
\[V_{dsat} \approx V_c\]
در این شرایط ولتاژ اشباع مستقل از ولتاژ گیت بوده و جریان اشباع به‌جای رابطه مربعی، رابطه‌ای خطی با ولتاژ گیت پیدا می‌کند.

\item \textbf{نتیجه کلی:}
در تمام ترانزیستورهای عملی نانومتری، مقدار \(V_{dsat}\) همواره از هر دو مقدار \(V_{GT}\) و \(V_c\) کوچک‌تر است:
\[V_{dsat} < \min(V_{GT}, V_c)\]
\end{enumerate}
